%!TeX root=../princesstop.tex
\chapter{The Visitor}

\begin{figure}[t!]
\centering
\includegraphics[width=\linewidth]{filigree}
\end{figure}

\lettrine[lines=5]{I}{magine}, if you can, what the rest of the evening was like. How they crouched by the fire which blazed and leaped and made so much of itself in the little grate. How they removed the covers of the dishes, and found rich, hot, savoury soup, which was a meal in itself, and sandwiches and toast and muffins enough for both of them. The mug from the washstand was used as Becky's tea cup, and the tea was so delicious that it was not necessary to pretend that it was anything but tea. They were warm and full-fed and happy, and it was just like Sara that, having found her strange good fortune real, she should give herself up to the enjoyment of it to the utmost. She had lived such a life of imaginings that she was quite equal to accepting any wonderful thing that happened, and almost to cease, in a short time, to find it bewildering.

<I don't know anyone in the world who could have done it,> she said; <but there has been someone. And here we are sitting by their fire—and—and—it's true! And whoever it is—wherever they are—I have a friend, Becky—someone is my friend.>

It cannot be denied that as they sat before the blazing fire, and ate the nourishing, comfortable food, they felt a kind of rapturous awe, and looked into each other's eyes with something like doubt.

<Do you think,> Becky faltered once, in a whisper, <do you think it could melt away, miss? Hadn't we better be quick?> And she hastily crammed her sandwich into her mouth. If it was only a dream, kitchen manners would be overlooked.

<No, it won't melt away,> said Sara. <I am \textsc{eating} this muffin, and I can taste it. You never really eat things in dreams. You only think you are going to eat them. Besides, I keep giving myself pinches; and I touched a hot piece of coal just now, on purpose.>

The sleepy comfort which at length almost overpowered them was a heavenly thing. It was the drowsiness of happy, well-fed childhood, and they sat in the fire glow and luxuriated in it until Sara found herself turning to look at her transformed bed.

There were even blankets enough to share with Becky. The narrow couch in the next attic was more comfortable that night than its occupant had ever dreamed that it could be.

As she went out of the room, Becky turned upon the threshold and looked about her with devouring eyes.

<If it ain't here in the mornin', miss,> she said, <it's been here tonight, anyways, an' I shan't never forget it.> She looked at each particular thing, as if to commit it to memory. <The fire was \textsc{there}>, pointing with her finger, <an' the table was before it; an' the lamp was there, an' the light looked rosy red; an' there was a satin cover on your bed, an' a warm rug on the floor, an' everythin' looked beautiful; an'>—she paused a second, and laid her hand on her stomach tenderly—<there \textsc{was} soup an' sandwiches an' muffins—there \textsc{was}.> And, with this conviction a reality at least, she went away.

Through the mysterious agency which works in schools and among servants, it was quite well known in the morning that Sara Crewe was in horrible disgrace, that Ermengarde was under punishment, and that Becky would have been packed out of the house before breakfast, but that a scullery maid could not be dispensed with at once. The servants knew that she was allowed to stay because Miss Minchin could not easily find another creature helpless and humble enough to work like a bounden slave for so few shillings a week. The elder girls in the schoolroom knew that if Miss Minchin did not send Sara away it was for practical reasons of her own.

<She's growing so fast and learning such a lot, somehow,> said Jessie to Lavinia, <that she will be given classes soon, and Miss Minchin knows she will have to work for nothing. It was rather nasty of you, Lavvy, to tell about her having fun in the garret. How did you find it out?>

<I got it out of Lottie. She's such a baby she didn't know she was telling me. There was nothing nasty at all in speaking to Miss Minchin. I felt it my duty>—priggishly. <She was being deceitful. And it's ridiculous that she should look so grand, and be made so much of, in her rags and tatters!>

<What were they doing when Miss Minchin caught them?>

<Pretending some silly thing. Ermengarde had taken up her hamper to share with Sara and Becky. She never invites us to share things. Not that I care, but it's rather vulgar of her to share with servant girls in attics. I wonder Miss Minchin didn't turn Sara out—even if she does want her for a teacher.>

<If she was turned out where would she go?> inquired Jessie, a trifle anxiously.

<How do I know?> snapped Lavinia. <She'll look rather queer when she comes into the schoolroom this morning, I should think—after what's happened. She had no dinner yesterday, and she's not to have any today.>

Jessie was not as ill-natured as she was silly. She picked up her book with a little jerk.

<Well, I think it's horrid,> she said. <They've no right to starve her to death.>

When Sara went into the kitchen that morning the cook looked askance at her, and so did the housemaids; but she passed them hurriedly. She had, in fact, overslept herself a little, and as Becky had done the same, neither had had time to see the other, and each had come downstairs in haste.

Sara went into the scullery. Becky was violently scrubbing a kettle, and was actually gurgling a little song in her throat. She looked up with a wildly elated face.

<It was there when I wakened, miss—the blanket,> she whispered excitedly. <It was as real as it was last night.>

<So was mine,> said Sara. <It is all there now—all of it. While I was dressing I ate some of the cold things we left.>

<Oh, laws! Oh, laws!> Becky uttered the exclamation in a sort of rapturous groan, and ducked her head over her kettle just in time, as the cook came in from the kitchen.

Miss Minchin had expected to see in Sara, when she appeared in the schoolroom, very much what Lavinia had expected to see. Sara had always been an annoying puzzle to her, because severity never made her cry or look frightened. When she was scolded she stood still and listened politely with a grave face; when she was punished she performed her extra tasks or went without her meals, making no complaint or outward sign of rebellion. The very fact that she never made an impudent answer seemed to Miss Minchin a kind of impudence in itself. But after yesterday's deprivation of meals, the violent scene of last night, the prospect of hunger today, she must surely have broken down. It would be strange indeed if she did not come downstairs with pale cheeks and red eyes and an unhappy, humbled face.

Miss Minchin saw her for the first time when she entered the schoolroom to hear the little French class recite its lessons and superintend its exercises. And she came in with a springing step, colour in her cheeks, and a smile hovering about the corners of her mouth. It was the most astonishing thing Miss Minchin had ever known. It gave her quite a shock. What was the child made of? What could such a thing mean? She called her at once to her desk.

<You do not look as if you realize that you are in disgrace,> she said. <Are you absolutely hardened?>

The truth is that when one is still a child—or even if one is grown up—and has been well fed, and has slept long and softly and warm; when one has gone to sleep in the midst of a fairy story, and has wakened to find it real, one cannot be unhappy or even look as if one were; and one could not, if one tried, keep a glow of joy out of one's eyes. Miss Minchin was almost struck dumb by the look of Sara's eyes when she made her perfectly respectful answer.

<I beg your pardon, Miss Minchin,> she said; <I know that I am in disgrace.>

<Be good enough not to forget it and look as if you had come into a fortune. It is an impertinence. And remember you are to have no food today.>

<Yes, Miss Minchin,> Sara answered; but as she turned away her heart leaped with the memory of what yesterday had been. <If the Magic had not saved me just in time,> she thought, <how horrible it would have been!>

<She can't be very hungry,> whispered Lavinia. <Just look at her. Perhaps she is pretending she has had a good breakfast>—with a spiteful laugh.

<She's different from other people,> said Jessie, watching Sara with her class. <Sometimes I'm a bit frightened of her.>

<Ridiculous thing!> ejaculated Lavinia.

All through the day the light was in Sara's face, and the colour in her cheek. The servants cast puzzled glances at her, and whispered to each other, and Miss Amelia's small blue eyes wore an expression of bewilderment. What such an audacious look of well-being, under august displeasure could mean she could not understand. It was, however, just like Sara's singular obstinate way. She was probably determined to brave the matter out.

One thing Sara had resolved upon, as she thought things over. The wonders which had happened must be kept a secret, if such a thing were possible. If Miss Minchin should choose to mount to the attic again, of course all would be discovered. But it did not seem likely that she would do so for some time at least, unless she was led by suspicion. Ermengarde and Lottie would be watched with such strictness that they would not dare to steal out of their beds again. Ermengarde could be told the story and trusted to keep it secret. If Lottie made any discoveries, she could be bound to secrecy also. Perhaps the Magic itself would help to hide its own marvels.

<But whatever happens,> Sara kept saying to herself all day—<\textsc{whatever} happens, somewhere in the world there is a heavenly kind person who is my friend—my friend. If I never know who it is—if I never can even thank him—I shall never feel quite so lonely. Oh, the Magic was \textsc{good} to me!>

If it was possible for weather to be worse than it had been the day before, it was worse this day—wetter, muddier, colder. There were more errands to be done, the cook was more irritable, and, knowing that Sara was in disgrace, she was more savage. But what does anything matter when one's Magic has just proved itself one's friend. Sara's supper of the night before had given her strength, she knew that she should sleep well and warmly, and, even though she had naturally begun to be hungry again before evening, she felt that she could bear it until breakfast-time on the following day, when her meals would surely be given to her again. It was quite late when she was at last allowed to go upstairs. She had been told to go into the schoolroom and study until ten o'clock, and she had become interested in her work, and remained over her books later.

When she reached the top flight of stairs and stood before the attic door, it must be confessed that her heart beat rather fast.

<Of course it \textsc{might} all have been taken away,> she whispered, trying to be brave. <It might only have been lent to me for just that one awful night. But it \textsc{was} lent to me—I had it. It was real.>

She pushed the door open and went in. Once inside, she gasped slightly, shut the door, and stood with her back against it looking from side to side.

The Magic had been there again. It actually had, and it had done even more than before. The fire was blazing, in lovely leaping flames, more merrily than ever. A number of new things had been brought into the attic which so altered the look of it that if she had not been past doubting she would have rubbed her eyes. Upon the low table another supper stood—this time with cups and plates for Becky as well as herself; a piece of bright, heavy, strange embroidery covered the battered mantel, and on it some ornaments had been placed. All the bare, ugly things which could be covered with draperies had been concealed and made to look quite pretty. Some odd materials of rich colours had been fastened against the wall with fine, sharp tacks—so sharp that they could be pressed into the wood and plaster without hammering. Some brilliant fans were pinned up, and there were several large cushions, big and substantial enough to use as seats. A wooden box was covered with a rug, and some cushions lay on it, so that it wore quite the air of a sofa.

Sara slowly moved away from the door and simply sat down and looked and looked again.

<It is exactly like something fairy come true,> she said. <There isn't the least difference. I feel as if I might wish for anything—diamonds or bags of gold—and they would appear! \textsc{that} wouldn't be any stranger than this. Is this my garret? Am I the same cold, ragged, damp Sara? And to think I used to pretend and pretend and wish there were fairies! The one thing I always wanted was to see a fairy story come true. I am \textsc{living} in a fairy story. I feel as if I might be a fairy myself, and able to turn things into anything else.>

She rose and knocked upon the wall for the prisoner in the next cell, and the prisoner came.

When she entered she almost dropped in a heap upon the floor. For a few seconds she quite lost her breath.

<Oh, laws!> she gasped. <Oh, laws, miss!>

<You see,> said Sara.

On this night Becky sat on a cushion upon the hearth rug and had a cup and saucer of her own.

When Sara went to bed she found that she had a new thick mattress and big downy pillows. Her old mattress and pillow had been removed to Becky's bedstead, and, consequently, with these additions Becky had been supplied with unheard-of comfort.

<Where does it all come from?> Becky broke forth once. <Laws, who does it, miss?>

<Don't let us even \textsc{ask},> said Sara. <If it were not that I want to say, <Oh, thank you,> I would rather not know. It makes it more beautiful.>

From that time life became more wonderful day by day. The fairy story continued. Almost every day something new was done. Some new comfort or ornament appeared each time Sara opened the door at night, until in a short time the attic was a beautiful little room full of all sorts of odd and luxurious things. The ugly walls were gradually entirely covered with pictures and draperies, ingenious pieces of folding furniture appeared, a bookshelf was hung up and filled with books, new comforts and conveniences appeared one by one, until there seemed nothing left to be desired. When Sara went downstairs in the morning, the remains of the supper were on the table; and when she returned to the attic in the evening, the magician had removed them and left another nice little meal. Miss Minchin was as harsh and insulting as ever, Miss Amelia as peevish, and the servants were as vulgar and rude. Sara was sent on errands in all weathers, and scolded and driven hither and thither; she was scarcely allowed to speak to Ermengarde and Lottie; Lavinia sneered at the increasing shabbiness of her clothes; and the other girls stared curiously at her when she appeared in the schoolroom. But what did it all matter while she was living in this wonderful mysterious story? It was more romantic and delightful than anything she had ever invented to comfort her starved young soul and save herself from despair. Sometimes, when she was scolded, she could scarcely keep from smiling.

<If you only knew!> she was saying to herself. <If you only knew!>

The comfort and happiness she enjoyed were making her stronger, and she had them always to look forward to. If she came home from her errands wet and tired and hungry, she knew she would soon be warm and well fed after she had climbed the stairs. During the hardest day she could occupy herself blissfully by thinking of what she should see when she opened the attic door, and wondering what new delight had been prepared for her. In a very short time she began to look less thin. colour came into her cheeks, and her eyes did not seem so much too big for her face.

<Sara Crewe looks wonderfully well,> Miss Minchin remarked disapprovingly to her sister.

<Yes,> answered poor, silly Miss Amelia. <She is absolutely fattening. She was beginning to look like a little starved crow.>

<Starved!> exclaimed Miss Minchin, angrily. <There was no reason why she should look starved. She always had plenty to eat!>

<Of—of course,> agreed Miss Amelia, humbly, alarmed to find that she had, as usual, said the wrong thing.

<There is something very disagreeable in seeing that sort of thing in a child of her age,> said Miss Minchin, with haughty vagueness.

<What—sort of thing?> Miss Amelia ventured.

<It might almost be called defiance,> answered Miss Minchin, feeling annoyed because she knew the thing she resented was nothing like defiance, and she did not know what other unpleasant term to use. <The spirit and will of any other child would have been entirely humbled and broken by—by the changes she has had to submit to. But, upon my word, she seems as little subdued as if—as if she were a princess.>

<Do you remember,> put in the unwise Miss Amelia, <what she said to you that day in the schoolroom about what you would do if you found out that she was\longdash>

<No, I don't,> said Miss Minchin. <Don't talk nonsense.> But she remembered very clearly indeed.

Very naturally, even Becky was beginning to look plumper and less frightened. She could not help it. She had her share in the secret fairy story, too. She had two mattresses, two pillows, plenty of bed-covering, and every night a hot supper and a seat on the cushions by the fire. The Bastille had melted away, the prisoners no longer existed. Two comforted children sat in the midst of delights. Sometimes Sara read aloud from her books, sometimes she learned her own lessons, sometimes she sat and looked into the fire and tried to imagine who her friend could be, and wished she could say to him some of the things in her heart.

Then it came about that another wonderful thing happened. A man came to the door and left several parcels. All were addressed in large letters, <To the Little Girl in the right-hand attic.>

Sara herself was sent to open the door and take them in. She laid the two largest parcels on the hall table, and was looking at the address, when Miss Minchin came down the stairs and saw her.

<Take the things to the young lady to whom they belong,> she said severely. <Don't stand there staring at them.>

<They belong to me,> answered Sara, quietly.

<To you?> exclaimed Miss Minchin. <What do you mean?>

<I don't know where they come from,> said Sara, <but they are addressed to me. I sleep in the right-hand attic. Becky has the other one.>

Miss Minchin came to her side and looked at the parcels with an excited expression.

<What is in them?> she demanded.

<I don't know,> replied Sara.

<Open them,> she ordered.

Sara did as she was told. When the packages were unfolded Miss Minchin's countenance wore suddenly a singular expression. What she saw was pretty and comfortable clothing—clothing of different kinds: shoes, stockings, and gloves, and a warm and beautiful coat. There were even a nice hat and an umbrella. They were all good and expensive things, and on the pocket of the coat was pinned a paper, on which were written these words: <To be worn every day. Will be replaced by others when necessary.>

Miss Minchin was quite agitated. This was an incident which suggested strange things to her sordid mind. Could it be that she had made a mistake, after all, and that the neglected child had some powerful though eccentric friend in the background—perhaps some previously unknown relation, who had suddenly traced her whereabouts, and chose to provide for her in this mysterious and fantastic way? Relations were sometimes very odd—particularly rich old bachelor uncles, who did not care for having children near them. A man of that sort might prefer to overlook his young relation's welfare at a distance. Such a person, however, would be sure to be crotchety and hot-tempered enough to be easily offended. It would not be very pleasant if there were such a one, and he should learn all the truth about the thin, shabby clothes, the scant food, and the hard work. She felt very queer indeed, and very uncertain, and she gave a side glance at Sara.

<Well,> she said, in a voice such as she had never used since the little girl lost her father, <someone is very kind to you. As the things have been sent, and you are to have new ones when they are worn out, you may as well go and put them on and look respectable. After you are dressed you may come downstairs and learn your lessons in the schoolroom. You need not go out on any more errands today.>

About half an hour afterward, when the schoolroom door opened and Sara walked in, the entire seminary was struck dumb.

<My word!> ejaculated Jessie, jogging Lavinia's elbow. <Look at the Princess Sara!>

Everybody was looking, and when Lavinia looked she turned quite red.

It was the Princess Sara indeed. At least, since the days when she had been a princess, Sara had never looked as she did now. She did not seem the Sara they had seen come down the back stairs a few hours ago. She was dressed in the kind of frock Lavinia had been used to envying her the possession of. It was deep and warm in colour, and beautifully made. Her slender feet looked as they had done when Jessie had admired them, and the hair, whose heavy locks had made her look rather like a Shetland pony when it fell loose about her small, odd face, was tied back with a ribbon.

<Perhaps someone has left her a fortune,> Jessie whispered. <I always thought something would happen to her. She's so queer.>

<Perhaps the diamond mines have suddenly appeared again,> said Lavinia, scathingly. <Don't please her by staring at her in that way, you silly thing.>

<Sara,> broke in Miss Minchin's deep voice, <come and sit here.>

And while the whole schoolroom stared and pushed with elbows, and scarcely made any effort to conceal its excited curiosity, Sara went to her old seat of honour, and bent her head over her books.

That night, when she went to her room, after she and Becky had eaten their supper she sat and looked at the fire seriously for a long time.

<Are you making something up in your head, miss?> Becky inquired with respectful softness. When Sara sat in silence and looked into the coals with dreaming eyes it generally meant that she was making a new story. But this time she was not, and she shook her head.

<No,> she answered. <I am wondering what I ought to do.>

Becky stared—still respectfully. She was filled with something approaching reverence for everything Sara did and said.

<I can't help thinking about my friend,> Sara explained. <If he wants to keep himself a secret, it would be rude to try and find out who he is. But I do so want him to know how thankful I am to him—and how happy he has made me. Anyone who is kind wants to know when people have been made happy. They care for that more than for being thanked. I wish—I do wish\longdash>

She stopped short because her eyes at that instant fell upon something standing on a table in a corner. It was something she had found in the room when she came up to it only two days before. It was a little writing-case fitted with paper and envelopes and pens and ink.

<Oh,> she exclaimed, <why did I not think of that before?>

She rose and went to the corner and brought the case back to the fire.

<I can write to him,> she said joyfully, <and leave it on the table. Then perhaps the person who takes the things away will take it, too. I won't ask him anything. He won't mind my thanking him, I feel sure.>

So she wrote a note. This is what she said:

\begin{quotation}
	
I hope you will not think it is impolite that I should write this note to you when you wish to keep yourself a secret. Please believe I do not mean to be impolite or try to find out anything at all; only I want to thank you for being so kind to me—so heavenly kind—and making everything like a fairy story. I am so grateful to you, and I am so happy—and so is Becky. Becky feels just as thankful as I do—it is all just as beautiful and wonderful to her as it is to me. We used to be so lonely and cold and hungry, and now—oh, just think what you have done for us! Please let me say just these words. It seems as if I \textsc{ought} to say them. \textsc{thank} you—\textsc{thank} you—\textsc{thank} you!

\begin{flushright}\scshape
The Little Girl In The Attic.
\end{flushright}
\end{quotation}


The next morning she left this on the little table, and in the evening it had been taken away with the other things; so she knew the Magician had received it, and she was happier for the thought. She was reading one of her new books to Becky just before they went to their respective beds, when her attention was attracted by a sound at the skylight. When she looked up from her page she saw that Becky had heard the sound also, as she had turned her head to look and was listening rather nervously.

<Something's there, miss,> she whispered.

<Yes,> said Sara, slowly. <It sounds—rather like a cat—trying to get in.>

She left her chair and went to the skylight. It was a queer little sound she heard—like a soft scratching. She suddenly remembered something and laughed. She remembered a quaint little intruder who had made his way into the attic once before. She had seen him that very afternoon, sitting disconsolately on a table before a window in the Indian gentleman's house.

<Suppose,> she whispered in pleased excitement—<just suppose it was the monkey who got away again. Oh, I wish it was!>

She climbed on a chair, very cautiously raised the skylight, and peeped out. It had been snowing all day, and on the snow, quite near her, crouched a tiny, shivering figure, whose small black face wrinkled itself piteously at sight of her.

<It is the monkey,> she cried out. <He has crept out of the Lascar's attic, and he saw the light.>

Becky ran to her side.

<Are you going to let him in, miss?> she said.

<Yes,> Sara answered joyfully. <It's too cold for monkeys to be out. They're delicate. I'll coax him in.>

She put a hand out delicately, speaking in a coaxing voice—as she spoke to the sparrows and to Melchisedec—as if she were some friendly little animal herself.

<Come along, monkey darling,> she said. <I won't hurt you.>

He knew she would not hurt him. He knew it before she laid her soft, caressing little paw on him and drew him towards her. He had felt human love in the slim brown hands of Ram Dass, and he felt it in hers. He let her lift him through the skylight, and when he found himself in her arms he cuddled up to her breast and looked up into her face.

<Nice monkey! Nice monkey!> she crooned, kissing his funny head. <Oh, I do love little animal things.>

He was evidently glad to get to the fire, and when she sat down and held him on her knee he looked from her to Becky with mingled interest and appreciation.

<He \textsc{is} plain-looking, miss, ain't he?> said Becky.

<He looks like a very ugly baby,> laughed Sara. <I beg your pardon, monkey; but I'm glad you are not a baby. Your mother \textsc{couldn't} be proud of you, and no one would dare to say you looked like any of your relations. Oh, I do like you!>

She leaned back in her chair and reflected.

<Perhaps he's sorry he's so ugly,> she said, <and it's always on his mind. I wonder if he \textsc{has} a mind. Monkey, my love, have you a mind?>

But the monkey only put up a tiny paw and scratched his head.

<What shall you do with him?> Becky asked.

<I shall let him sleep with me tonight, and then take him back to the Indian gentleman tomorrow. I am sorry to take you back, monkey; but you must go. You ought to be fondest of your own family; and I'm not a \textsc{real} relation.>

And when she went to bed she made him a nest at her feet, and he curled up and slept there as if he were a baby and much pleased with his quarters.